\begin{exsol}{6}{Tadelis Exercise 4.5}{77}{Foundations}{B3}
\asked Find the strategy profiles of the given $3 \times 3$ game that survive IESDS.

\method IESDS round by round (Lecture~1, slides~28--29; Tadelis \S4.2.1, p.~65): in each round
remove every strategy that is strictly dominated given the opponent's surviving strategies.

\solution
Matrix as printed in the book (p.~77):
\begin{center}
\begin{tabular}{r|c|c|c|}
\multicolumn{1}{r}{1 $\backslash$ 2} & \multicolumn{1}{c}{$L$} & \multicolumn{1}{c}{$C$} & \multicolumn{1}{c}{$R$}\\ \cline{2-4}
$U$ & $6,\ 8$ & $2,\ 6$ & $8,\ 2$ \\ \cline{2-4}
$M$ & $8,\ 2$ & $4,\ 4$ & $9,\ 5$ \\ \cline{2-4}
$D$ & $8,\ 10$ & $4,\ 6$ & $6,\ 7$ \\ \cline{2-4}
\end{tabular}
\end{center}

\textbf{Round 1} ($S^0$ = full game).
\begin{itemize}
  \item Player~1: $U$ is strictly dominated by $M$: $6 < 8$, $2 < 4$, $8 < 9$. $D$ is not
  dominated by $M$ (ties $8 = 8$ and $4 = 4$), and $M$, $D$ are not dominated by $U$.
  \item Player~2: her payoffs are $L = (8, 2, 10)$, $C = (6, 4, 6)$, $R = (2, 5, 7)$ in rows
  $(U, M, D)$. $C$ is beaten by $L$ in rows $U, D$ but not in $M$; it beats $R$ in row $U$. No
  column is strictly dominated.
\end{itemize}
$S_1^1 = \{M, D\}$, $S_2^1 = \{L, C, R\}$.

\textbf{Round 2.} Player~2 now compares columns on rows $(M, D)$ only:
$L = (2, 10)$, $C = (4, 6)$, $R = (5, 7)$. $C$ is strictly dominated by $R$ ($4 < 5$, $6 < 7$);
$L$ and $R$ are not ranked. Player~1 on all three columns: $M = (8, 4, 9)$, $D = (8, 4, 6)$:
$M$ is better only against $R$, so $D$ is weakly but not strictly dominated.
$S_1^2 = \{M, D\}$, $S_2^2 = \{L, R\}$.

\textbf{Round 3.} Player~1 on $(L, R)$: $M = (8, 9)$, $D = (8, 6)$, a tie against $L$, so no
strict dominance. Player~2 on $(M, D)$: $L = (2, 10)$, $R = (5, 7)$, not ranked. The process
stops.
\[
  S^\infty = \{M, D\} \times \{L, R\} = \{(M,L), (M,R), (D,L), (D,R)\}.
\]

\interp Rationality of player~1 removes $U$. Then player~2's knowledge that player~1 is rational
removes $C$, which can be a best reply only to beliefs that put positive weight on $U$. The iterated-elimination
equilibrium is a set of four profiles, not a point. Two of them, $(M,R)$ with payoffs $(9,5)$
and $(D,L)$ with $(8,10)$, are profiles where each strategy is a best reply to the other (the
pure Nash equilibria of Lecture~2). Deleting the weakly dominated $D$ as well would single out
$(M,R)$, which is a stronger and less safe assumption (see item~19).

\verify \texttt{check\_A.py} runs IESDS with the lecture's definition and finds exactly the
rounds above (round~1: $U$; round~2: $C$; then stop), that $U$ is dominated only by $M$ and $C$
only by $R$, and that $D$ is weakly dominated by $M$ in the final set.

\answer{Round~1: delete $U$ (dominated by $M$). Round~2: delete $C$ (dominated by $R$ once $U$
is gone). Nothing else is strictly dominated:
$S^\infty = \{M, D\} \times \{L, R\}$, four profiles.}
\end{exsol}
